% Modern editorial, markup and figure/layout contributions: CC-BY-SA-4.0.
% Historical text and illustrations: Leonhard Euler (1744), public domain.
% Component scope and upstream attribution: see RIGHTS.md.
\subsection*{Hypothesis I.}
\originalsection{16}\textit{In hac tractatione abscissam, ad quam omnes curvas referemus, perpetuo littera $x$, applicatam vero littera $y$ designabimus. Tum vero, sumptis elementis abscissæ æqualibus, semper erit}
\[
 dy=pdx;\quad dp=qdx;\quad dq=rdx;\quad dr=sdx;\quad\text{\&c.}
\]
\subsection*{Coroll. I.}
\originalsection{17}His igitur substitutionibus omnia differentialia ipsius $y$ cujuscunque gradus ex expressionibus tollentur, atque præter differentiale $dx$ nulla alia differentialia relinquentur. Quanquam autem hoc modo omnia differentialia, præter $dx$, specie tantum, non revera tolluntur; tamen hæ substitutiones ingens nobis in præsenti instituto afferent subsidium.
\subsection*{Coroll. II.}
\originalsection{18}Quin etiam hujusmodi substitutionibus differentialis constantis assumtio penitus de calculo tollitur: quodcunque enim differentiale aliud constans assumatur, post istas substitutiones perpetuo eadem formula emergere debet. Interim tamen, ob methodum infra adhibendam, necesse erit differentiale $dx$ tanquam constans assumere.
\sourcepage{9}
\subsection*{Coroll. III.}
\originalsection{19}Ut autem facilius appareat, quomodo per has substitutiones differentialia cujusque gradus ipsius $y$ evanescant; juvabit sequentem Tabellam adjecisse.
\[
\begin{aligned}
 dy&=pdx\\
 ddy&=dpdx=qdx^2\\
 d^3y&=dqdx^2=rdx^3\\
 d^4y&=drdx^3=sdx^4\\
 d^5y&=dsdx^4=tdx^5\\
 &\text{\&c.}\qquad\text{\&c.}\qquad\text{\&c.}
\end{aligned}
\]
\subsection*{Coroll. IV.}
\originalsection{20}Quod si etiam arcus curvæ abscissæ $x$ respondens, cum suis differentialibus cujuscunque gradus occurrat; ea omnia per istas litteras ita exprimi poterunt, ut nulla alia differentialia præter $dx$ adsint. Posito enim arcu $=w$ erit.
\[
\begin{aligned}
 w&=\int\sqrt{(dx^2+dy^2)}=\int dx\sqrt{(1+pp)}\\
 dw&=dx\sqrt{(1+pp)}\\
 ddw&=\frac{pqdx^2}{\sqrt{(1+pp)}}\\
 d^3w&=\frac{prdx^3}{\sqrt{(1+pp)}}+\frac{qqdx^3}{(1+pp)^{3:2}}\\
 &\text{\&c.}
\end{aligned}
\]
\subsection*{Coroll. V.}
\originalsection{21}Simili modo, ex his radius osculi seu curvedinis curvæ, in quovis loco, per quantitates specie saltem finitas poterit exprimi. Cum enim, posito elemento $dx$ constante, sit longitudo radii osculi $=\dfrac{-dw^3}{dxddy}$; fiet ea $=\dfrac{-(1+pp)^{3:2}}{q}$.
